\documentclass[10pt,a4paper]{amsart}
\usepackage[margin=19mm]{geometry}
\usepackage{amsmath,amssymb,mathtools}
\usepackage[scr=rsfs,cal=euler]{mathalfa}
\usepackage{xfrac}
\usepackage[unicode,hidelinks,bookmarks=false]{hyperref}
\newcommand{\ie}{i.e.\ }
\newcommand{\cf}{cf.\ }
\newcommand{\resp}{respectively\ }
\newcommand{\Subsection}{\subsection}
\providecommand{\nobreakdash}{\nobreak\hbox{-}}
\setlength{\emergencystretch}{2em}
\multlinegap0pt
\usepackage{amssymb}
\newtheorem{lemma}{Lemma}
\title{Separation for Schr\"odinger operators. A counterexample to Simon's conjecture}
\author{M.A. Perelmuter}
\date{Source: arXiv:2609.00415v1 (CC BY 4.0)}
\begin{document}
\maketitle

\noindent\emph{Source and licence.} Separation for Schr\"odinger operators. A counterexample to Simon's conjecture. Version arXiv:2609.00415v1 (CC BY 4.0), distributed by arXiv under the Creative Commons Attribution 4.0 International licence, http://creativecommons.org/licenses/by/4.0/, as recorded in the arXiv licence metadata for this exact version and shown on the abstract page https://arxiv.org/abs/2609.00415v1. Full licence notice: \texttt{licenses/arxiv-2609.00415-CC-BY-4.0.txt}.\par\smallskip

\noindent\textit{Editorial scope.} Counted unit: the article's lemma lem:gluing (source0.tex lines 201--206). The article's complete original proof of it is delivered with it (source0.tex lines 207--225, one \texttt{proof} environment of 19 lines). No mathematics is written, repaired or re-derived: the statement and the proof are the article's own lines, in the article's own notation, and no retained line is substituted or re-flowed.  No further source-local context is needed: every cross-reference of every retained line resolves inside the two delivered spans.  Nothing was omitted from any retained span, so the package carries no omission record.  No retained line contains a citation, so the wrapper carries no bibliography and no bibliography entry.  No retained line contains a figure environment, an image-inclusion command or drawing code, so no image is delivered and the package declares no asset dependency.  Editorial formatting only, in addition: an AMS \texttt{amsart} preamble that restates the article's own declarations and macros used by the retained lines, namely the environments \texttt{lemma} on separate counters, together with the standard packages those lines need.  The retained environments print in this document as Lemma 1 in order of appearance, whereas the article prints the same items with the numbering of its own full document.  The complete native source of the article as distributed in the arXiv e-print for this exact version is bundled unchanged as \texttt{sources/209024/source0.tex}.
\begin{lemma}\label{lem:gluing}
Let $d\ge 3$, and let $u$ be the solution of
$$-\left(r^{d-1}u'(r)\right)'=r^{d-1},\qquad 1\le r\le 2,$$
subject to the initial conditions $u(1)=cU(1)$, $u'(1)=cU'(1)$, where $c,\,U(1),\,U'(1)$ are given real constants. Then
$$u(2)+\frac{2}{d-2}u'(2)=c\left(U(1)+\frac{U'(1)}{d-2}\right)-\frac{3}{2(d-2)}.$$
\end{lemma}

\begin{proof}
The general solution of the inhomogeneous equation is
$$u(r)=A+Br^{2-d}-\frac{r^2}{2d} \quad \text{and} \quad u'(r)=B(2-d)r^{1-d}-\frac{r}{d}.$$
Imposing the initial conditions at $r=1$:
$$u(1)=A+B-\frac{1}{2d}=cU(1),\qquad u'(1)=B(2-d)-\frac1d=cU'(1).$$
These equations give
$$B=-\frac{cU'(1)+\frac1d}{d-2}=-\frac{cU'(1)}{d-2}-\frac{1}{d(d-2)},\quad A=cU(1)-B+\frac{1}{2d}.$$
Now evaluate at $r=2$:
$$u(2)=A+B\,2^{2-d}-\frac2d,\qquad u'(2)=B(2-d)2^{1-d}-\frac2d.$$
The second equality gives
$$\frac{2}{d-2}u'(2)=\frac{2}{d-2}\Big[B(2-d)2^{1-d}-\frac2d\Big]=-B\,2^{2-d}-\frac{4}{d(d-2)}.$$
Adding this to $u(2)$:
$$u(2)+\frac{2}{d-2}u'(2)=A-\frac2d-\frac{4}{d(d-2)}.$$
Substituting $A=cU(1)-B+\dfrac{1}{2d}$ and then $B=-\dfrac{cU'(1)}{d-2}-\dfrac{1}{d(d-2)}$:
$$u(2)+\frac{2}{d-2}u'(2)=cU(1)+\frac{cU'(1)}{d-2}+\frac{1}{d(d-2)}-\frac{3}{2d}-\frac{4}{d(d-2)}=c\left(U(1)+\frac{U'(1)}{d-2}\right)-\frac{3}{d(d-2)}-\frac{3}{2d}.$$
Finally, the remaining constants combine as
$$\frac{3}{d(d-2)}+\frac{3}{2d}=\frac{3}{2(d-2)},$$
which yields the claimed identity.
\end{proof}

\end{document}
